% Source: https://artofproblemsolving.com/community/c3h2170056_easy_inequality
% Stable identity: aops-topic-2170056-q01
Call the left-hand side $E$. Putting $c=a+b$ makes $E$ vanish (try it), which suggests the factor $a+b-c$, and by the symmetry of $E$ in $b$ and $c$ also $a-b+c$. Indeed,
\[(a+b-c)(a-b+c)=a^2-(b-c)^2=a^2-b^2+2bc-c^2,\]
and multiplying by $a+b+c$,
\begin{align*}
(a+b+c)(a^2-b^2-c^2+2bc)&=\left(a^3-ab^2-ac^2+2abc\right)+\left(a^2b-b^3-bc^2+2b^2c\right)+\left(a^2c-b^2c-c^3+2bc^2\right)\\
&=a^3+a^2b+a^2c-ab^2+2abc-ac^2-b^3+b^2c+bc^2-c^3=E .
\end{align*}
So
\[E=(a+b+c)(a+b-c)(a-b+c).\]

\begin{asy}
usepackage("amsmath");
size(4cm);
defaultpen(fontsize(10pt));
pair A=(0,0), B=(5,0), C=(1.6,3);
draw(A--B--C--cycle,linewidth(0.9));
label("$c$",(A+B)/2,S); label("$a$",(B+C)/2,NE); label("$b$",(A+C)/2,NW);
label("$A$",A,SW); label("$B$",B,SE); label("$C$",C,N);
label("\mbox{}",(min(currentpicture,user=true)+max(currentpicture,user=true))/2,black); // keeps the text after the figure black when LaTeX typesets the labels inline
\end{asy}

In a triangle with sides $a$, $b$, $c$, the triangle inequality gives $a+b>c$ and $a+c>b$, and of course $a+b+c>0$. All three factors are positive, so $E>0$; in particular $E\ge0$. (Equality would need a degenerate triangle with $c=a+b$ or $b=a+c$.) $\blacksquare$

\emph{Remark.} With the substitution $a=y+z$, $b=z+x$, $c=x+y$ ($x,y,z>0$), which describes every triangle, one gets $E=8yz(x+y+z)$.
